% 03-diffusion
\chapter{Diffusion}\label{ch:diffusion}

Before we let reactions and diffusion compete, we must understand what diffusion
does on its own. This chapter is our first encounter with a partial differential
equation, the diffusion equation $u_t=d\,u_{xx}$. We derive it, add the zero-flux
boundary condition used throughout these notes, and show in two ways that diffusion
makes every profile uniform: by an energy estimate, and by expanding in cosines. The
second method turns the PDE into decoupled scalar ODEs, one for each cosine mode; it
is the prototype of everything we do with two species in \cref{ch:turing}.

\section{Where the diffusion equation comes from}\label{sec:diffusion-derivation}

Consider a substance in a thin tube, modeled by the interval $[0,L]$. We assume that
the concentration varies only along the tube and denote by $u(x,t)$ the amount of
substance per unit length at position $x$ and time $t$. Let the \emph{flux}
$J(x,t)$ be the net amount per unit time that passes the point $x$
from left to right, and let $f(x,t)$ be the net rate of production by reactions per
unit length ($f<0$ means consumption). The amount in a piece $[x,x+h]$ changes by what
enters at $x$, minus what leaves at $x+h$, plus what is produced inside:
\begin{equation*}
  \frac{\mathrm d}{\mathrm dt}\int_x^{x+h}u(y,t)\dd y
  = J(x,t)-J(x+h,t)+\int_x^{x+h}f(y,t)\dd y .
\end{equation*}
For smooth $u$, $J$, $f$ we differentiate under the integral sign, divide by $h$ and
let $h\to0$. This gives the \emph{conservation law}
\begin{equation}\label{eq:diffusion-conservation}
  u_t = -J_x + f ,
\end{equation}
valid whatever the mechanism of transport. For diffusion, i.e.\ transport by the
random thermal motion of molecules, the flux obeys \emph{Fick's law}
\begin{equation}\label{eq:fick}
  J = -d\,u_x
\end{equation}
with a constant \emph{diffusion coefficient} $d>0$: substance flows from high to low
concentration, at a rate proportional to the gradient. Inserting \eqref{eq:fick} into
\eqref{eq:diffusion-conservation} gives, without reactions and with reactions that
depend only on the local concentration,
\begin{align}
  u_t &= d\,u_{xx} &&\text{(diffusion equation)},\label{eq:heat}\\
  u_t &= d\,u_{xx}+f(u) &&\text{(reaction--diffusion equation)}.\label{eq:diffusion-reaction}
\end{align}
Since \eqref{eq:heat} also describes heat conduction, it is also called the
\emph{heat equation}.

\paragraph{Units and the diffusive time scale.}
As $u_t$ has the units of $u$ per second and $u_{xx}$ those of $u$ per square meter,
$d$ is measured in $\mathrm{m}^2/\mathrm{s}$, and for every length $\ell$ the
quantity $\ell^2/d$ is a time: the \emph{diffusive time scale} on which diffusion
evens out differences over the distance $\ell$; doubling $\ell$ quadruples it.
(Intuitively, a molecule performing a random walk is after $n$ steps typically
only about $\sqrt n$ steps away from its start.) For small molecules in water,
$d\approx10^{-9}\,\mathrm{m}^2/\mathrm{s}$; this gives $0.1\,\mathrm{s}$ for
$\ell=10\,\mu\mathrm{m}$ (a cell), about $17$ minutes for $1\,\mathrm{mm}$, and more
than a day for $1\,\mathrm{cm}$.

\section{Boundary and initial conditions}\label{sec:diffusion-bc}

Throughout these notes the tube is \emph{closed}: nothing enters or leaves through
its ends, $J(0,t)=J(L,t)=0$. By Fick's law this is the homogeneous \emph{Neumann} or
\emph{zero-flux} boundary condition
\begin{equation}\label{eq:neumann}
  u_x(0,t)=u_x(L,t)=0,\qquad t>0 .
\end{equation}
It says nothing about the values of $u$ at the ends. No spatial structure is imposed
from outside, so any pattern that forms under \eqref{eq:neumann} is produced inside
the domain: the natural setting for pattern formation. Other conditions appear only
in the exercises: the \emph{Dirichlet} condition $u(0,t)=u(L,t)=0$ (ends in contact
with a reservoir that removes the substance) and \emph{periodic} conditions
$u(0,t)=u(L,t)$, $u_x(0,t)=u_x(L,t)$ (a ring of length $L$). With an initial profile $u_0$ we arrive at the problem
\begin{equation}\label{eq:heat-problem}
\begin{aligned}
  u_t &= d\,u_{xx}, & x\in(0,L),\ t>0,\\
  u_x(0,t) &= u_x(L,t) = 0, & t>0,\\
  u(x,0) &= u_0(x), & x\in(0,L).
\end{aligned}
\end{equation}

\begin{remark}[Standing assumption]\label{rem:well-posedness}
Throughout these notes we assume that for smooth initial data compatible with the
boundary conditions (for \eqref{eq:heat-problem}: $u_0'(0)=u_0'(L)=0$) the problems
we consider have exactly one solution, and that it is \emph{smooth}: $u$ and $u_x$
are continuous on $[0,L]\times[0,\infty)$, and all partial derivatives of $u$ exist
and are continuous on $[0,L]\times(0,\infty)$.
For nonlinear reactions, solutions may exist only for a finite time (for
$u_t=u_{xx}+u^2$ and positive constant initial data they blow up); statements then
refer to the interval of existence. Proofs belong to a PDE course
\cite{Evans2010,Strauss2008}.
\end{remark}

\section{Mass conservation}\label{sec:diffusion-mass}

\begin{proposition}[Conservation of mass]\label{prop:mass-conservation}
For a smooth solution $u$ of \eqref{eq:heat-problem},
$\frac{\mathrm d}{\mathrm dt}\int_0^Lu(x,t)\dd x=0$ for $t>0$. Hence
$\int_0^Lu(x,t)\dd x=\int_0^Lu_0\dd x$ for all $t\ge0$: the mean value of $u(\cdot,t)$
is constant in time and equal to $M=\frac1L\int_0^Lu_0\dd x$.
\end{proposition}

\begin{proof}
For $t>0$ we may differentiate under the integral sign, since $u_t$ is continuous.
By the equation, the fundamental theorem of calculus and \eqref{eq:neumann},
\begin{equation*}
  \frac{\mathrm d}{\mathrm dt}\int_0^L u\dd x = \int_0^L u_t\dd x
  = d\int_0^L u_{xx}\dd x = d\,\big(u_x(L,t)-u_x(0,t)\big) = 0 .
\end{equation*}
So $m(t)=\int_0^Lu(x,t)\dd x$ is constant on $(0,\infty)$. As $u$ is uniformly
continuous on $[0,L]\times[0,1]$, $m$ is continuous at $t=0$, so $m(t)=m(0)$.
\end{proof}

For \eqref{eq:diffusion-reaction} the same computation gives
$\frac{\mathrm d}{\mathrm dt}\int_0^Lu\dd x=\int_0^Lf(u)\dd x$: with zero flux,
the total amount changes only through reactions.

\section{Diffusion makes things uniform}\label{sec:diffusion-energy}

If $u(\cdot,t)$ becomes constant as $t\to\infty$, the constant must be $M$. We prove
that this happens, exponentially fast, in the norm
$\norm{u}=\big(\int_0^Lu^2\dd x\big)^{1/2}$; we write $u(t)$ for $x\mapsto u(x,t)$.

Let $u$ be a smooth solution of \eqref{eq:heat-problem}. Constants solve
\eqref{eq:heat} and satisfy \eqref{eq:neumann}, so $w=u-M$ does too. For $t>0$ we
multiply $w_t=d\,w_{xx}$ by $w$, integrate, and integrate by parts:
\begin{equation}\label{eq:diffusion-energy}
  \frac{\mathrm d}{\mathrm dt}\norm{u-M}^2
  = 2\int_0^L w\,w_t\dd x
  = 2d\big[w\,w_x\big]_0^L - 2d\int_0^L w_x^2\dd x
  = -2d\norm{u_x}^2 .
\end{equation}
The boundary term vanishes because $w_x=u_x=0$ at $x=0,L$. So $\norm{u(t)-M}$ never
increases. (For the difference $w=u_1-u_2$ of two smooth solutions with the same
initial datum, the same computation shows that $\norm{w(t)}$ is nonincreasing, and
it vanishes at $t=0$; so \eqref{eq:heat-problem} has at most one smooth solution.)
To obtain a rate, we compare $\norm{u_x}$ with $\norm{u-M}$: a function with mean
value zero cannot be large unless its derivative is large.

\begin{lemma}[Poincar\'e--Wirtinger inequality]\label{lem:poincare-wirtinger}
Let $u\in C^1[0,L]$ with $\int_0^Lu\dd x=0$. Then
$\norm{u}^2 \le (L/\pi)^2\norm{u'}^2$, and the constant $(L/\pi)^2$ is the best
possible.
\end{lemma}

\begin{proof}
We use Fourier series (motivated in \cref{sec:diffusion-separation}, treated in
\cref{app:fourier}). A square-integrable $g$ on $(0,L)$ has the \emph{cosine
coefficients} $\hat g_0=\frac1L\int_0^Lg\dd x$,
$\hat g_k=\frac2L\int_0^Lg(x)\cos(k\pi x/L)\dd x$ and the \emph{sine coefficients}
$b_k=\frac2L\int_0^Lg(x)\sin(k\pi x/L)\dd x$ ($k\ge1$), and it is the sum of its
cosine series $\sum_{k\ge0}\hat g_k\cos(k\pi x/L)$ and of its sine series
$\sum_{k\ge1}b_k\sin(k\pi x/L)$. Since these functions are mutually orthogonal,
Pythagoras' theorem takes the form of \emph{Parseval's identity}
$\norm{g}^2=L\hat g_0^2+\frac L2\sum_{k\ge1}\hat g_k^2=\frac L2\sum_{k\ge1}b_k^2$
(\cref{thm:cosine-series}). Let $\hat u_k$ be the cosine coefficients of $u$ and
$b_k$ the sine coefficients of $u'$. Integrating by parts,
\begin{equation}\label{eq:diffusion-sine-coefficients}
  b_k = \frac2L\Big[u(x)\sin\frac{k\pi x}{L}\Big]_0^L
        - \frac{k\pi}{L}\,\frac2L\int_0^L u(x)\cos\frac{k\pi x}{L}\dd x
      = -\frac{k\pi}{L}\,\hat u_k ,\qquad k\ge1 .
\end{equation}
The boundary term vanishes because $\sin(0)=\sin(k\pi)=0$; note that \emph{no
boundary condition on $u$} is needed. By Parseval's identity for $u'$ and for $u$,
and since $\hat u_0=0$,
\begin{equation*}
  \norm{u'}^2 = \frac L2\sum_{k=1}^\infty b_k^2
  = \frac L2\sum_{k=1}^\infty \Big(\frac{k\pi}{L}\Big)^2\hat u_k^2
  \ge \Big(\frac{\pi}{L}\Big)^2\,\frac L2\sum_{k=1}^\infty \hat u_k^2
  = \Big(\frac{\pi}{L}\Big)^2\norm{u}^2 .
\end{equation*}
For $u(x)=\cos(\pi x/L)$ we have $\int_0^Lu\dd x=0$, $\norm{u}^2=L/2$ and
$\norm{u'}^2=(\pi/L)^2L/2$; so equality holds, and the constant cannot be improved.
\end{proof}

\begin{theorem}[Diffusion homogenizes]\label{thm:heat-decay}
Let $u$ be a smooth solution of \eqref{eq:heat-problem} and
$M=\frac1L\int_0^Lu_0\dd x$. Then
\begin{equation*}
  \norm{u(t)-M} \le e^{-d\pi^2t/L^2}\,\norm{u_0-M}\qquad\text{for all } t\ge0 .
\end{equation*}
\end{theorem}

\begin{proof}
Let $E(t)=\norm{u(t)-M}^2$ and $\gamma=d\pi^2/L^2$. For $t>0$, $u(t)-M$ lies in
$C^1[0,L]$ and has mean value zero (\cref{prop:mass-conservation}), so
\cref{lem:poincare-wirtinger} gives $\norm{u_x(t)}^2\ge(\pi/L)^2E(t)$, and
\eqref{eq:diffusion-energy} becomes $E'\le-2\gamma E$. Hence
$\frac{\mathrm d}{\mathrm dt}\big(e^{2\gamma t}E\big)=e^{2\gamma t}(E'+2\gamma E)\le0$
on $(0,\infty)$. Since $E$ is continuous on $[0,\infty)$ (as in the proof of
\cref{prop:mass-conservation}), $E(t)\le e^{-2\gamma t}E(0)$ for all $t\ge0$.
\end{proof}

\begin{keyidea}
With zero-flux boundary conditions, diffusion \emph{homogenizes}: every solution of
the heat equation converges to its mean value $M$, exponentially fast with rate at
least $d\pi^2/L^2$. Its time scale $(L/\pi)^2/d$ is the diffusive time scale of the
domain: small domains and fast diffusion homogenize faster.
\end{keyidea}

\section{Separation of variables}\label{sec:diffusion-separation}

The energy method tells us \emph{that} $u$ approaches $M$, not \emph{how} the profile
changes. To see this, we look for solutions $u(x,t)=T(t)\varphi(x)$ with
$\varphi\not\equiv0$. Then \eqref{eq:heat} reads $T'\varphi=d\,T\varphi''$, and
wherever $T\varphi\neq0$ we get $T'(t)/(d\,T(t))=\varphi''(x)/\varphi(x)$. The
left-hand side does not depend on $x$, the right-hand side not on $t$, so both equal a
constant $-\lambda$. Thus $T(t)=e^{-d\lambda t}T(0)$, and, by \eqref{eq:neumann},
$\varphi$ solves the \emph{Neumann eigenvalue problem}
\begin{equation}\label{eq:neumann-eigenproblem}
  -\varphi'' = \lambda\varphi\quad\text{on } (0,L),\qquad \varphi'(0)=\varphi'(L)=0 .
\end{equation}
We call $\lambda\in\R$ an \emph{eigenvalue} if \eqref{eq:neumann-eigenproblem} has
a solution $\varphi\in C^2[0,L]$, $\varphi\not\equiv0$ (an \emph{eigenfunction}).

\begin{proposition}[Neumann eigenpairs of the interval]\label{prop:neumann-eigen}
The eigenvalues of \eqref{eq:neumann-eigenproblem} are exactly
\begin{equation*}
  \lambda_k = \Big(\frac{k\pi}{L}\Big)^2,\qquad k=0,1,2,\dots,
\end{equation*}
and the eigenfunctions for $\lambda_k$ are the nonzero multiples of
$\varphi_k(x)=\cos(k\pi x/L)$. In particular, $\lambda_0=0$ is an eigenvalue, with
the constant eigenfunction $\varphi_0\equiv1$.
\end{proposition}

\begin{proof}
Let $\varphi$ solve \eqref{eq:neumann-eigenproblem}; we use the general solution of
$\varphi''+\lambda\varphi=0$. \emph{Case $\lambda=-\sigma^2<0$, $\sigma>0$.} Then
$\varphi=\alpha\cosh(\sigma x)+\beta\sinh(\sigma x)$, so
$\varphi'(0)=\sigma\beta=0$ gives $\beta=0$, and
$\varphi'(L)=\sigma\alpha\sinh(\sigma L)=0$ gives $\alpha=0$ since
$\sinh(\sigma L)>0$. So $\varphi\equiv0$: no eigenvalue.

\emph{Case $\lambda=0$.} Then $\varphi=\alpha+\beta x$ and $\varphi'(0)=\beta=0$, so
$\varphi$ is constant.

\emph{Case $\lambda=\sigma^2>0$, $\sigma>0$.} Then
$\varphi=\alpha\cos(\sigma x)+\beta\sin(\sigma x)$, so $\varphi'(0)=\sigma\beta=0$
gives $\beta=0$, and $\varphi'(L)=-\sigma\alpha\sin(\sigma L)=0$. A solution with
$\alpha\neq0$ exists if and only if $\sin(\sigma L)=0$, i.e.\ $\sigma L=k\pi$ with
$k\in\N$; then $\lambda=(k\pi/L)^2$ and $\varphi=\alpha\cos(k\pi x/L)$. Conversely,
$\varphi_k''=-\lambda_k\varphi_k$, and $\varphi_k'(x)=-\frac{k\pi}{L}\sin(k\pi x/L)$
vanishes at $x=0$ and $x=L$.
\end{proof}

The eigenfunctions (\cref{fig:ch03-heat-modes}(a)) are mutually orthogonal:
\begin{equation}\label{eq:diffusion-orthogonality}
  \int_0^L\varphi_k\varphi_m\dd x = 0\quad (k\neq m),\qquad
  \int_0^L\varphi_k^2\dd x = \frac L2\quad (k\ge1),\qquad
  \int_0^L\varphi_0^2\dd x = L .
\end{equation}
Indeed, $\cos\alpha\cos\beta=\frac12\big(\cos(\alpha-\beta)+\cos(\alpha+\beta)\big)$
and $\int_0^L\cos(j\pi x/L)\dd x=0$ for every integer $j\neq0$; for $k\neq m$ both
$k\pm m$ are such integers, and for $k=m\ge1$ only the constant $\frac12$ survives.
If $u=\sum_k\hat u_k\varphi_k$, multiplying by $\varphi_m$ and integrating gives
$\int_0^Lu\varphi_m\dd x=\hat u_m\int_0^L\varphi_m^2\dd x$. This motivates the
\emph{cosine series} and the \emph{cosine coefficients} of a function $u$:
\begin{equation}\label{eq:cosine-coefficients}
  u(x) = \sum_{k=0}^\infty\hat u_k\cos\Big(\frac{k\pi x}{L}\Big),\quad
  \hat u_0 = \frac1L\int_0^L u\dd x,\quad
  \hat u_k = \frac2L\int_0^L u(x)\cos\Big(\frac{k\pi x}{L}\Big)\dd x\ \ (k\ge1).
\end{equation}
So $\hat u_0$ is the mean value of $u$. By \cref{thm:cosine-series} the series
converges to $u$ in mean square, and uniformly if $u\in C^1[0,L]$. In cosine
coefficients, the second derivative with Neumann
conditions acts like a diagonal matrix:

\begin{lemma}[Cosine coefficients of the second derivative]\label{lem:coefficient-uxx}
Let $u\in C^2[0,L]$ with $u'(0)=u'(L)=0$. Then the cosine coefficients of $u''$ are
$-\lambda_k\hat u_k$, $k=0,1,2,\dots$, where $\hat u_k$ are those of $u$.
\end{lemma}

\begin{proof}
Let $k\ge1$. Integrating by parts twice,
\begin{equation*}
  \int_0^L u''\varphi_k\dd x
  = \big[u'\varphi_k\big]_0^L - \int_0^L u'\varphi_k'\dd x
  = \big[u'\varphi_k\big]_0^L - \big[u\,\varphi_k'\big]_0^L + \int_0^L u\,\varphi_k''\dd x .
\end{equation*}
The first boundary term vanishes because $u'(0)=u'(L)=0$: this is where the Neumann
condition on $u$ is used. The second vanishes because
$\varphi_k'(x)=-\frac{k\pi}{L}\sin(k\pi x/L)$ is zero at $x=0$ and $x=L$: this is a
property of the eigenfunction and holds for every $u$. Since
$\varphi_k''=-\lambda_k\varphi_k$, we get
$\frac2L\int_0^Lu''\varphi_k\dd x=-\lambda_k\hat u_k$. For $k=0$,
$\frac1L\int_0^Lu''\dd x=\frac1L\big(u'(L)-u'(0)\big)=0=-\lambda_0\hat u_0$.
\end{proof}

\begin{theorem}[The heat equation in cosine modes]\label{thm:heat-modes}
Let $u$ be a smooth solution of \eqref{eq:heat-problem}, and let $\hat u_k(t)$ be the
cosine coefficients of $u(t)$ (so $\hat u_k(0)$ are those of $u_0$). Then for every
$k\ge0$
\begin{equation*}
  \hat u_k'(t) = -d\lambda_k\,\hat u_k(t),\qquad\text{hence}\qquad
  \hat u_k(t) = e^{-d\lambda_kt}\,\hat u_k(0).
\end{equation*}
Consequently, for every $t\ge0$, with uniform convergence in $x\in[0,L]$,
\begin{equation}\label{eq:heat-series}
  u(x,t) = \sum_{k=0}^\infty e^{-d\lambda_kt}\,\hat u_k(0)\cos\Big(\frac{k\pi x}{L}\Big).
\end{equation}
\end{theorem}

\begin{proof}
Let $k\ge1$ and $t>0$. Differentiating under the integral sign and using the equation,
\begin{equation*}
  \hat u_k'(t) = \frac2L\int_0^L u_t(x,t)\,\varphi_k(x)\dd x
  = d\,\frac2L\int_0^L u_{xx}(x,t)\,\varphi_k(x)\dd x = -d\lambda_k\,\hat u_k(t),
\end{equation*}
by \cref{lem:coefficient-uxx} applied to $u(t)$, which satisfies the Neumann
condition. For $k=0$ the same computation (with $\frac1L$ in place of $\frac2L$)
gives $\hat u_0'=0$, as in \cref{prop:mass-conservation}. So
$\frac{\mathrm d}{\mathrm dt}\big(e^{d\lambda_kt}\hat u_k(t)\big)=0$ for $t>0$, and
since $\hat u_k$ is continuous on $[0,\infty)$, $\hat u_k(t)=e^{-d\lambda_kt}\hat u_k(0)$.
Finally, $u(t)\in C^1[0,L]$, so its cosine series converges uniformly to $u(t)$
(\cref{thm:cosine-series}); this is \eqref{eq:heat-series}.
\end{proof}

Each term of \eqref{eq:heat-series} is a separated solution
$e^{-d\lambda_kt}\varphi_k(x)$, and every solution is a superposition of them. The
series makes sense for every square-integrable $u_0$, even a discontinuous one, and
still defines a smooth solution for $t>0$ \cite{Strauss2008}; this is how
\cref{fig:ch03-heat-modes}(b) was computed.

\begin{keyidea}
In cosine coefficients the heat equation decouples into the scalar ODEs
$\hat u_k'=-d\lambda_k\hat u_k$, $k=0,1,2,\dots$, one for each mode. Mode $k$ decays
at the rate $d\lambda_k=d\,k^2\pi^2/L^2$: short wavelengths die first, and only the
mean $\hat u_0=M$ is preserved.
\end{keyidea}

\paragraph{A second proof of \cref{thm:heat-decay}.}
Since $\hat u_0(t)=M$, the function $u(t)-M$ has the cosine coefficients
$0,\hat u_1(t),\hat u_2(t),\dots$, and Parseval's identity \eqref{eq:parseval} gives
\begin{equation*}
  \norm{u(t)-M}^2 = \frac L2\sum_{k=1}^\infty e^{-2d\lambda_kt}\,\hat u_k(0)^2
  \le e^{-2d\lambda_1t}\,\frac L2\sum_{k=1}^\infty\hat u_k(0)^2
  = e^{-2d\lambda_1t}\,\norm{u_0-M}^2,
\end{equation*}
because $\lambda_k\ge\lambda_1=\pi^2/L^2$ for $k\ge1$. The rate is sharp: for
$u_0=M+\cos(\pi x/L)$ the solution is $u=M+e^{-d\lambda_1t}\cos(\pi x/L)$, with
equality.

The rate $d\lambda_k$ grows quadratically in $k$, so at time $t$ only modes with
$d\lambda_kt\lesssim1$ are visible: structures on scales below about $\sqrt{dt}$
are erased. In \cref{fig:ch03-heat-modes}(b), small-scale noise disappears almost
instantly, while the step is smoothed out much more slowly.

\begin{figure}[t]
  \centering
  \includegraphics[width=\textwidth]{ch03-heat-modes}
  \caption{(a) The Neumann eigenfunctions $\varphi_k(x)=\cos(k\pi x/L)$,
  $k=0,\dots,3$; $\varphi_k$ has $k$ zeros in $(0,L)$. (b) The solution of
  \eqref{eq:heat-problem} with $d=1$, $L=1$, computed from the cosine series
  \eqref{eq:heat-series}. The initial profile (grey) is a step plus noise made of the
  modes $k=12,\dots,40$. The noise is gone by $t=10^{-2}$ (damping factor at most
  $e^{-d\lambda_{12}t}<10^{-6}$), while the step is only smoothed out. For large $t$
  the profile approaches the mean value $M=0.48$ (dashed) in the shape of $\varphi_1$.}
  \label{fig:ch03-heat-modes}
\end{figure}

\section{With reactions}\label{sec:diffusion-reactions}

Diffusion alone homogenizes, and stable kinetics alone return to their steady state.
The question for the rest of these notes is: \emph{can the combination be
destabilizing?} For one species, \eqref{eq:diffusion-reaction} with the Neumann
condition, each zero $\ub$ of $f$ is a constant steady state. Linearizing about
$\ub$ and expanding in cosines as in \cref{thm:heat-modes}, mode $k$ obeys a scalar
ODE with rate $f'(\ub)-d\lambda_k$, so diffusion seems only to stabilize;
\cref{ch:one-equation} makes this precise. For two species, $u_t=d_u\,u_{xx}+f(u,v)$,
$v_t=d_v\,v_{xx}+g(u,v)$ (\cref{ch:turing}), each cosine mode of the linearized
problem is a $2\times2$ linear ODE system, and the answer can be yes:
\emph{diffusion-driven instability (DDI)} or \emph{Turing instability}
\cite{Turing1952}.

\section{Beyond the interval}\label{sec:diffusion-general}

Animal coat patterns live in two dimensions. On a bounded domain $\Omega\subset\R^n$
the diffusion equation is $u_t=d\,\Delta u$, $\Delta u=\sum_{i}\partial^2u/\partial x_i^2$,
with the zero-flux condition $\partial_\nu u=\nabla u\cdot\nu=0$ on $\partial\Omega$
($\nu$ the outer unit normal). Mass conservation and the energy identity carry over
(via the divergence theorem), and separation of variables leads to
$-\Delta\varphi=\lambda\varphi$ in $\Omega$, $\partial_\nu\varphi=0$ on
$\partial\Omega$, whose solutions behave like cosines.

\begin{theorem}[Neumann eigenpairs of a bounded domain]\label{thm:neumann-eigen-general}
Let $\Omega\subset\R^n$ be a bounded, connected open set with smooth boundary. There
are numbers $0=\lambda_0<\lambda_1\le\lambda_2\le\dots\to\infty$ and functions
$\varphi_0,\varphi_1,\dots$, smooth on $\overline\Omega$, with
$-\Delta\varphi_k=\lambda_k\varphi_k$ in $\Omega$ and $\partial_\nu\varphi_k=0$ on
$\partial\Omega$, that form an orthonormal basis of $L^2(\Omega)$. Each eigenvalue is
repeated according to the dimension of its eigenspace, and $\varphi_0$ is constant.
\end{theorem}

We omit the proof; see \cite[Ch.~6]{Evans2010}, where the Dirichlet case is treated
and the Neumann case is analogous. ``Orthonormal basis'' means that
$\int_\Omega\varphi_j\varphi_k\dd x$ is $1$ for $j=k$ and $0$ otherwise, and that
every square-integrable $u$ is the mean-square limit of $\sum_{k\le N}c_k\varphi_k$,
$c_k=\int_\Omega u\varphi_k\dd x$, as for cosine series; on $(0,L)$ the normalized
cosines $1/\sqrt L$, $\sqrt{2/L}\cos(k\pi x/L)$ form such a basis. As on the
interval, indexing starts at $k=0$; by connectedness, only constants have eigenvalue
$0$.

\begin{example}[Rectangle]\label{ex:rectangle-eigen}
On $\Omega=(0,L_1)\times(0,L_2)$ the products
\begin{equation*}
  \varphi_{k,l}(x,y)=\cos\Big(\frac{k\pi x}{L_1}\Big)\cos\Big(\frac{l\pi y}{L_2}\Big),
  \qquad
  \lambda_{k,l}=\pi^2\Big(\frac{k^2}{L_1^2}+\frac{l^2}{L_2^2}\Big),
  \qquad k,l=0,1,2,\dots,
\end{equation*}
satisfy $-\Delta\varphi_{k,l}=\lambda_{k,l}\varphi_{k,l}$ and
$\partial_\nu\varphi_{k,l}=0$ (on the sides $x=0,L_1$, $\partial_\nu=\mp\partial_x$
produces the factor $\sin(k\pi x/L_1)=0$; similarly for $y=0,L_2$), and, normalized,
they form an orthonormal basis (by cosine series in each variable). The modes $\varphi_{k,0}$ and $\varphi_{0,l}$
depend on one variable only: they are \emph{stripes}. For $k,l\ge1$ the sign of
$\varphi_{k,l}$ alternates like a checkerboard, and the maxima form a lattice of
\emph{spots}. Several modes can share an eigenvalue: on a square, $L_1=L_2=L$, the
stripes $\varphi_{1,0}$, $\varphi_{0,1}$ and all their combinations have the
eigenvalue $\pi^2/L^2$ (so $\lambda_1=\lambda_2$), and $5^2+0^2=4^2+3^2$ gives a
four-dimensional eigenspace for $25\pi^2/L^2$ containing stripes and spots. Which of
them a pattern-forming system selects is beyond linear analysis (\cref{ch:examples}).
\end{example}

\begin{remark}[Wavenumber]\label{rem:wavenumber}
We call $q_k=\sqrt{\lambda_k}=k\pi/L$ the \emph{wavenumber} of mode $k$. Then
$\varphi_k(x)=\cos(q_kx)$ has the \emph{wavelength} $2\pi/q_k=2L/k$: mode $k$ fits
$k$ half-wavelengths into $(0,L)$. In higher dimensions $\lambda=\abs{\mathbf q}^2$
for a wave vector $\mathbf q$, e.g.\ $\mathbf q=(k\pi/L_1,l\pi/L_2)$ on the rectangle.
Some authors call $\lambda_k$ itself the ``wave number''; we do not.
\end{remark}

\section{Exercises}\label{sec:diffusion-exercises}

\begin{exercise}[Dirichlet conditions]\label{exer:diffusion-dirichlet}
(a) Show that the eigenvalues of $-\varphi''=\lambda\varphi$ on $(0,L)$ with
$\varphi(0)=\varphi(L)=0$ are exactly $(k\pi/L)^2$, $k\in\N$, with eigenfunctions
$\sin(k\pi x/L)$ (up to factors). (b) Let $u$ be a smooth solution of
$u_t=d\,u_{xx}$ with $u(0,t)=u(L,t)=0$ and $u(x,0)=u_0(x)$. Show that the sine
coefficients $b_k(t)$ of $u(t)$ satisfy $b_k'=-d(k\pi/L)^2b_k$, and deduce
$\norm{u(t)}\le e^{-d\pi^2t/L^2}\norm{u_0}$. (c) Is mass conserved? Compute
$\frac{\mathrm d}{\mathrm dt}\int_0^Lu\dd x$ and interpret it.
\end{exercise}

\begin{exercise}[Periodic conditions]\label{exer:diffusion-periodic}
Show that the eigenvalues of $-\varphi''=\lambda\varphi$ on $(0,L)$ with
$\varphi(0)=\varphi(L)$, $\varphi'(0)=\varphi'(L)$ are $(2k\pi/L)^2$,
$k=0,1,2,\dots$; the eigenfunctions are the constants for $k=0$ and the
combinations of $\cos(2k\pi x/L)$, $\sin(2k\pi x/L)$ for $k\ge1$. Is mass conserved?
\end{exercise}

\begin{exercise}\label{exer:diffusion-cos2}
Solve \eqref{eq:heat-problem} explicitly for $u_0(x)=\cos^2(\pi x/L)$. Determine $M$
and show that $\norm{u(t)-M}=e^{-4d\pi^2t/L^2}\norm{u_0-M}$. Why is the decay faster
than guaranteed by \cref{thm:heat-decay}? Which symmetry of $u_0$ explains it?
\end{exercise}

\begin{exercise}[Sharpness of Poincar\'e--Wirtinger]\label{exer:diffusion-poincare}
(a) Show that equality holds in \cref{lem:poincare-wirtinger} if and only if
$u(x)=c\cos(\pi x/L)$, $c\in\R$. (b) Show that the inequality fails without the
assumption $\int_0^Lu\dd x=0$, but that $\norm{u-\hat u_0}\le(L/\pi)\norm{u'}$ for
every $u\in C^1[0,L]$.
\end{exercise}

\begin{exercise}[$\ast$]\label{exer:diffusion-max}
Let $u$ be a smooth solution of \eqref{eq:heat-problem}. (a) Show that
$\frac{\mathrm d}{\mathrm dt}\int_0^Lu^4\dd x=-12d\int_0^Lu^2u_x^2\dd x\le0$ for $t>0$,
and more generally $\frac{\mathrm d}{\mathrm dt}\int_0^L\Phi(u)\dd x\le0$ for every
convex $\Phi\in C^2(\R)$. (b) Deduce that
$\max_x\abs{u(x,t)}\le\max_x\abs{u_0(x)}$ for all $t\ge0$. \emph{Hint:} take
$\Phi(s)=s^{2p}$ and use $\big(\int_0^L\abs{g}^{2p}\dd x\big)^{1/(2p)}\to\max\abs{g}$
as $p\to\infty$ for continuous $g$.
\end{exercise}
